\documentclass[conference]{IEEEtran}
\IEEEoverridecommandlockouts
\usepackage{cite}
\usepackage{amsmath,amssymb,amsfonts}
\usepackage{graphicx}
\usepackage{textcomp}
\usepackage{xcolor}
\usepackage{booktabs}
\usepackage{url}

% Registered results resolve through docs/research/result-traceability.md
% and docs/research/findings.md. The preliminary encoder comparison uses
% runs/encoders-20261022/checkpoints/stream-lipsync-avhubert-history.json.

\begin{document}

\title{Multimodal Audio-Visual Deepfake Detection:\\
Validity Checks Catch a Detector That Looks Right and Is Not}

\author{\IEEEauthorblockN{Vaibhav Shiroorkar}
\IEEEauthorblockA{\textit{Dept. of Electronics and Computer Engineering} \\
\textit{KJ Somaiya School of Engineering}\\
Mumbai, India \\
v.shiroorkar@somaiya.edu}
\and
\IEEEauthorblockN{Adithya Jithesh}
\IEEEauthorblockA{\textit{Dept. of Electronics and Computer Engineering} \\
\textit{KJ Somaiya School of Engineering}\\
Mumbai, India \\
adithya.j@somaiya.edu}
\and
\IEEEauthorblockN{Shivansh Gupta}
\IEEEauthorblockA{\textit{Dept. of Electronics and Computer Engineering} \\
\textit{KJ Somaiya School of Engineering}\\
Mumbai, India \\
shivansh.gupta@somaiya.edu}
\and
\IEEEauthorblockN{Tejashree Phatak}
\IEEEauthorblockA{\textit{Dept. of Electronics and Computer Engineering} \\
\textit{KJ Somaiya School of Engineering}\\
Mumbai, India \\
tejashree@somaiya.edu}
}

\maketitle

\begin{abstract}
Lip-sync deepfakes manipulate the mouth region and the accompanying voice
while leaving most of the original video untouched. This creates a problem
for detectors that rely only on visual information, since much of the footage
they analyse may still be genuine. Combining audio and video seems like a
natural solution, which led us to develop a system with separate visual and
audio streams and two cross-attention streams designed to compare mouth
movements with speech and facial information with voice characteristics.
Their outputs are combined at the feature level. Rather than assuming that
these components work as intended, however, we examined whether their
behaviour supports that assumption. We asked whether cue-specific fusion
could generalize better than a strong visual baseline under source-disjoint
evaluation and controlled testing for shortcut learning. Our results did not
support this expectation. Fusion improved the in-domain ROC-AUC by 0.0159,
but it failed to outperform the visual baseline on an unseen corpus across
all three seeds and performed significantly worse for two of them. The
experiments also exposed problems that would have been easy to overlook if
we had considered accuracy alone. One cross-attention stream achieved a
ROC-AUC of 0.9991 while its attention weights remained almost uniform
throughout training, suggesting that the mechanism it was designed around
was not contributing as intended. We also found evidence of codec-related
shortcut learning: when genuine clips were re-encoded using the codec
associated with the largest fake category, 93.3 percent were classified as
fake. Combining compression augmentation with codec-swapped checkpoint
selection reduced cross-corpus ROC-AUC by 0.19. In another experiment,
selecting a checkpoint using validation loss resulted in below-chance
performance, while changing only the training seed produced a cross-corpus
ROC-AUC spread of 0.1078. These findings show why it is important to
monitor properties that are not directly optimized during training. We
present our evaluation protocol, the measurements collected, and the
failures identified throughout the study.
\end{abstract}

\begin{IEEEkeywords}
Audio-visual deepfakes, cross-attention, cross-corpus generalization,
feature-level fusion, media forensics, shortcut learning.
\end{IEEEkeywords}

\section{Introduction}

Early deepfake techniques mainly focused on replacing entire faces. More
recent approaches can produce much more targeted manipulations. For example,
an attacker can generate speech using text-to-speech or voice conversion and
then regenerate only the mouth region to match the new audio. The rest of the
video, including its background, lighting, and upper facial region, can remain
unchanged. This creates a challenge for detectors that rely entirely on
visual information, since most of the footage they examine may still be
genuine.

One possible way to address this problem is to analyse audio and video
together and look for inconsistencies between them. These may include mouth
movements that do not match the spoken sounds or differences between the
emotion conveyed by a voice and the expression visible on the face. Our
system follows this approach. Visual streams process sequences of face crops,
an audio stream analyses speech, and two cross-attention streams compare the
mouth with the audio and the face with the voice. A learned head then combines
these representations at the feature level.

Before running the experiments, we defined the main research question and
allowed for the possibility that multimodal fusion might not improve
performance. Specifically, we investigated whether combining visual
artifacts, audio spoofing cues, and mouth-audio alignment could generalize
better than a strong visual baseline under source-disjoint evaluation with
controls for shortcut learning. The results did not support this expectation,
and we report them together with confidence intervals.

The more useful outcome of the study was identifying issues that would have
been difficult to notice from performance tables alone. Four findings were
particularly important:

\begin{enumerate}
\item One cross-attention stream achieved near-perfect performance even
though its attention weights remained close to their uniform starting values
throughout training.
\item The visual models learned to some extent from video codec differences,
because the largest fake category in one dataset used a different codec from
the genuine clips.
\item Our checkpoint selection procedure retained a model that performed
below chance during evaluation.
\item The training seed had a substantial effect on cross-corpus performance,
sometimes exceeding the differences observed between architectural choices.
\end{enumerate}

Each of these issues became apparent through a quantity that we recorded
alongside the training objective but did not directly optimize. We consider
this practice, rather than the architecture alone, to be the main contribution
of our work. It encourages us to check whether a model behaves as its design
suggests, instead of relying only on the final prediction scores. We also
apply the same approach to fully generated video, where a frozen encoder and
a linear probe are evaluated on generators that were not seen during
training.

\section{Related Work}

Several datasets have been developed to study audio-visual deepfakes,
including FakeAVCeleb~\cite{fakeavceleb} and LAV-DF~\cite{lavdf}, which
contain manipulated faces and voices. Cross-dataset evaluation commonly
involves training on FaceForensics++~\cite{ffpp} and testing on datasets such
as Celeb-DF~\cite{celebdf} and DFDC~\cite{dfdc}. The problem of synchronizing
mouth movements with speech has also been studied in SyncNet~\cite{syncnet},
while AV-HuBERT~\cite{avhubert} learns a shared audio-visual speech
representation from unlabelled video. Self-blended images~\cite{sbi} take a
different approach by training face forgery detectors on blending artifacts
rather than on a particular generator. Similarly, a frozen vision-language
encoder combined with a linear probe has been shown to transfer across image
generators~\cite{univfd}.

Our work is also related to shortcut learning~\cite{shortcut}, where a model
achieves good benchmark results by relying on information that was not
intended to be part of the task. A concurrent study auditing deepfake
benchmarks using frozen self-supervised features~\cite{audit} is particularly
relevant to our findings on the choice of region of interest.

We do not claim that multimodal fusion itself is a new contribution. Instead,
we focus on a problem that can arise when evaluating such systems: a stream
may produce impressive scores even when the mechanism it was designed to use
does not appear to be functioning. We investigate this problem through
specific experiments and checks on the behaviour of the models.

\section{Method}

\subsection{Two systems, one evaluation path}

We investigate two system designs. Design A contains three branches. The
first uses EfficientNet-B0~\cite{efficientnet}, followed by a GRU that
processes 16 face crops of size $224\times224$. The crops are detected and
tracked using MTCNN~\cite{mtcnn}. The second branch uses Wav2Vec2~\cite{wav2vec2}
with attentive pooling over 4~s of 16~kHz audio to detect spoofing. The third
branch combines a framewise ResNet-18 mouth encoder with a second Wav2Vec2
model. It processes a 2~s window containing 50 mouth crops to estimate
mouth-audio alignment.

The three branches are combined using regularized logistic regression. Their
logits are calibrated using Platt scaling, and three quality features are
included in the fusion head. The fusion model is fitted only on out-of-fold
rows, keeping the fusion training separate from the predictions used to fit
it.

Design B uses configurable streams and replaces late fusion with a
feature-level fusion head. Each visual stream uses either EfficientNet-B0 or
a frozen DINOv3 ViT-S/16~\cite{dinov3} backbone followed by a temporal model.
The audio-visual streams use cross-attention to compare temporal tokens. The
attention operation is defined as

\begin{equation}
A = \operatorname{softmax}\!\left(\frac{(X_a W_Q)(X_v W_K)^\top}{\sqrt{d}}\right),
\qquad \tilde{X} = A\,(X_v W_V),
\end{equation}

where the audio tokens $X_a$ act as queries, while the video tokens $X_v$
provide the keys and values. The stream then embeds the residual
$X_a-\tilde{X}$ and averages it over time. The intended interpretation is
that this residual captures audio information that cannot be explained by the
corresponding mouth movements. The lip-sync stream uses the mouth window and
its associated audio, whereas the emotion stream uses face crops and voice
features.

For feature-level fusion, each stream produces a 256-dimensional embedding,
which is projected to a shared width before the embeddings are concatenated.
Missing streams are masked after projection, ensuring that an absent stream
contributes exactly zero. A multi-layer perceptron then produces the final
classification.

Both designs use the same evaluation code. When we first observed the largest
performance gap, we suspected that it might have resulted from an evaluation
bug. To check this, we evaluated the Design A checkpoint through the Design B
pipeline. It reproduced the same score to four decimal places, suggesting
that the difference was not caused by the evaluation path.

\subsection{A validity check recorded and never optimized}

A stream designed to align mouth movements with speech should assign greater
attention to temporally corresponding audio and video tokens. In particular,
we expect its attention matrix $A$ to place more weight near the diagonal,
allowing for a small offset between the mouth movement and the recorded sound.
We therefore record the diagonal mass after every epoch without including it
in the training loss:

\begin{equation}
m = \frac{1}{T_q}\sum_{i,j} A_{ij}\,\mathbb{1}\!\left[\,\left|\tfrac{T_k}{T_q}\,i - j\right| \le 1\right]
\end{equation}

The reason for keeping this measurement separate from the objective is that
optimizing a diagnostic measure can make it less useful as an independent
check of the mechanism. If the model is explicitly trained to maximize the
measure, the result no longer provides the same evidence about whether the
mechanism has emerged naturally during learning.

The expected value under uniform attention depends on the query length. For
50 steps, the expected value is 0.0592, while for 8 steps it is 0.3438.
These values provide a reference for determining whether the measured
attention is meaningfully different from a uniform distribution.

\subsection{Abstention and provenance}

Some clips do not contain a stable face track. Rather than falling back to
the full frame or silently excluding these clips, the pipeline marks the
visual view as unavailable. We report coverage alongside the evaluation
metrics so that the effect of missing views remains visible.

The data splits are fixed by source identity. Cached views are identified
using a hash of the preprocessing settings and code version, and the model
refuses to use a cache produced under different settings. Each registered
result is also stored with SHA-256 hashes of its artifacts, making it possible
to trace the reported measurements to the corresponding files.

The preliminary encoder comparison reported later in this paper is based on
a training history. It has not yet been evaluated on held-out data, so those
results should be interpreted as preliminary rather than as evidence of
cross-corpus performance.

\section{Experimental Setup}

For the first track, we train on FakeAVCeleb~\cite{fakeavceleb} using a
70/15/15 split based on source identity. DFDC~\cite{dfdc} serves as the
cross-corpus test set and is the only corpus with audio in this study that is
not derived from VoxCeleb2. We use Celeb-DF-v2~\cite{celebdf} and the
evaluation-only MNW benchmark~\cite{mnw} for additional external checks.
LAV-DF~\cite{lavdf} provides a stricter setting for evaluating the cross-modal
streams, while the c23 version of FaceForensics++~\cite{ffpp} is used as a
second training corpus following the published cross-dataset protocol.

The second track uses DF26~\cite{df26}, which contains 271 real clips and
2,420 generated clips from seven generators. Four of these generators are
commercial and three are open source. We train a frozen-encoder linear probe
on six generators and evaluate it on the remaining unseen generator.

Our evaluation protocol specifies three training seeds: 17, 29, and 43.
Every Design A branch and the fusion head were trained using all three seeds.
However, the Design B streams were trained only with seed 17, which happens
to give the best DFDC result among the three seeds in the broader comparison.
Consequently, the single-seed cross-corpus results for Design B should be
viewed as a best-case observation rather than a typical estimate of
performance.

We calculate 95 percent bootstrap confidence intervals~\cite{bootstrap},
with resampling clustered by source identity. This reflects the practical
question of how a detector might behave when presented with a new set of
speakers. If a resample contains only one class, we skip it and count it
rather than assigning it a score of 0.5. Comparisons between systems are
paired: both systems are evaluated on the same resampled identities, and
their score difference is calculated within each resample. For DF26, where
generated clips are not associated with speakers in the same way, we resample
clips or source videos instead.

\section{Results}

\subsection{The cross-modal streams do not read correspondence}

On held-out FakeAVCeleb clips, the emotion stream achieves a ROC-AUC of
0.9991, the highest score among the five streams. However, its attention
weights remain close to their uniform value throughout training. Table
\ref{tab:diag} compares the measured diagonal attention mass for the emotion
and lip-sync streams with their expected uniform values across the recorded
training epochs.

\begin{table}[htbp]
\caption{Diagonal attention mass against its uniform value}
\label{tab:diag}
\centering
\begin{tabular}{lrrl}
\toprule
Stream & Steps & Uniform & Measured \\
\midrule
emotion & 8 & 0.3438 & 0.3434 to 0.3457, 9 epochs \\
lip-sync & 50 & 0.0592 & 0.0592 to 0.0595, 4 epochs \\
lip-sync, lower encoder LR & 50 & 0.0592 & 0.0592 to 0.0595, 6 epochs \\
\bottomrule
\end{tabular}
\end{table}

Over nine epochs, the emotion stream's ROC-AUC increased from 0.9847 to
0.9991, but its attention distribution barely changed. This suggests that the
stream learned to distinguish the classes without relying on the temporal
correspondence that its design was intended to capture. The performance is
real, but the attention measurements do not support the intended explanation
for that performance.

FakeAVCeleb contains clips in which the manipulations can affect the video
as a whole. A model can therefore distinguish genuine and fake clips by
recognizing generator-related artifacts without comparing the mouth movements
with the audio. LAV-DF allows us to examine the stream under a stricter
condition. We extract two windows from each file: one containing the
manipulated segment and another that does not. Both windows share the same
speaker, lighting, microphone, and codec.

A lip-sync stream trained on these paired windows achieves a ROC-AUC of
0.9969 on 1,161 held-out windows. Its diagonal attention mass is 0.0592,
compared with a uniform value of 0.0600. Although matching the windows removes
many file-level differences, local artifacts introduced by the manipulation
can still distinguish the two classes. The near-uniform attention therefore
does not establish that the model has learned temporal correspondence.

\subsection{The visual models partly read the video codec}

We found a substantial codec difference in FakeAVCeleb. Genuine videos are
encoded in H.264 at approximately 150~kbps, whereas every fake in the
face-swap plus Wav2Lip category, the dataset's largest fake category, is
encoded in MPEG-4 Part~2 at approximately 330~kbps. To investigate whether
this difference affected the predictions, we re-encoded 75 test clips for
each condition and evaluated the deployed ensemble using a threshold of 0.5.
The results are shown in Table~\ref{tab:codec}.

\begin{table}[htbp]
\caption{Called fake after re-encoding (75 clips per row)}
\label{tab:codec}
\centering
\begin{tabular}{llr}
\toprule
Clips & Codec & Called fake \\
\midrule
genuine & original H.264 & 0.0\% \\
genuine & re-encoded H.264 (control) & 0.0\% \\
genuine & re-encoded MPEG-4 & \textbf{93.3\%} \\
swap plus Wav2Lip & original MPEG-4 & 100.0\% \\
swap plus Wav2Lip & re-encoded H.264 & \textbf{30.7\%} \\
face swap, H.264 & re-encoded MPEG-4 & 100.0\% \\
\bottomrule
\end{tabular}
\end{table}

The control condition helps separate the effect of re-encoding from the
effect of changing the codec. Re-encoding genuine clips with their original
codec changes the encoded video, but none of the clips are classified as fake.
Changing the codec produces a very different outcome: 93.3 percent of genuine
clips become false positives, while re-encoding the swap-plus-Wav2Lip clips
with H.264 reduces the fake detection rate to 30.7 percent. These changes
indicate that the predictions are sensitive to codec-related information.

Since decoding produces the pixels generated by the encoder, including its
compression artifacts, a model can learn to use those artifacts as a shortcut.
The results do not imply that every prediction is based on the codec, but they
show that codec information can have a substantial effect on the deployed
ensemble.

\subsection{Training against the codec cue costs transfer}

We compared four EfficientNet-B0 training recipes using seed 17. Each recipe
changes one component relative to the preceding one. The first is the shipped
recipe, which selects checkpoints using ROC-AUC. The second adds random JPEG,
H.264, and MPEG-4 round trips to both classes and selects checkpoints using
codec-swapped validation clips. The third adds self-blended fakes~\cite{sbi},
and the fourth includes the FaceForensics++ training split.

\begin{table}[htbp]
\caption{Visual recipes against the codec cue (ROC-AUC)}
\label{tab:recipes}
\centering
\begin{tabular}{lrrr}
\toprule
Recipe & FakeAVCeleb & DFDC & Celeb-DF \\
\midrule
shipped & 0.9988 & \textbf{0.7611} & 0.6682 \\
+ compression & 0.9959 & 0.5720 & 0.6525 \\
+ self-blending & 0.9978 & 0.6496 & 0.6762 \\
+ FaceForensics++ & 0.9904 & 0.6756 & 0.6509 \\
\bottomrule
\end{tabular}
\end{table}

As Table~\ref{tab:recipes} shows, all three modified recipes perform worse
than the shipped recipe on DFDC. The largest decrease comes from the
compression recipe, which loses 0.1891 ROC-AUC, and the corresponding
confidence intervals do not overlap. Self-blending and the additional
training corpus recover part of this loss, but their DFDC scores remain
0.0855 below the shipped recipe.

The results on Celeb-DF are less conclusive. The confidence intervals overlap
for all four recipes, so these experiments do not provide a clear basis for
distinguishing their performance on that corpus.

It is also important to separate the effects of the individual changes. The
compression augmentation and the codec-swapped checkpoint selection rule were
introduced together, so the experiment cannot establish how much of the
transfer loss comes specifically from reducing the model's reliance on codec
information. We also retrained the shipped recipe and reproduced the original
checkpoint exactly, epoch by epoch, confirming that the observed results were
repeatable under that training configuration.

\subsection{Fusion does not beat the visual branch across corpora}

Table~\ref{tab:fusion} shows the paired ROC-AUC difference between the fused
system and a model fitted using only the visual branch and the same quality
features. Each seed uses its own cross-fitted fold models.

\begin{table}[htbp]
\caption{Fusion minus visual, paired, 95\% interval}
\label{tab:fusion}
\centering
\begin{tabular}{lll}
\toprule
Seed & FakeAVCeleb test & DFDC \\
\midrule
17 & $+0.0159$ [$+0.0129$, $+0.0190$] & $-0.0085$ [$-0.0196$, $+0.0027$] \\
29 & $+0.0170$ [$+0.0143$, $+0.0195$] & $-0.0215$ [$-0.0371$, $-0.0052$] \\
43 & $+0.0184$ [$+0.0159$, $+0.0209$] & $-0.0417$ [$-0.0698$, $-0.0129$] \\
\bottomrule
\end{tabular}
\end{table}

Fusion improves performance on the in-domain test set for all three seeds.
Almost all of this improvement comes from 75 clips in which the video is
genuine but the voice has been replaced. A visual-only branch cannot identify
these cases from the video alone.

The pattern changes on DFDC. Fusion does not outperform the visual branch
for any of the three seeds and performs significantly worse for seeds 29
and 43. We also evaluated every subset of the five Design B streams using
the same rows. The best combination on DFDC, consisting of the emotion and
lip-sync streams, achieves 0.6076 ROC-AUC. Combining all five streams gives
only 0.5411, which is lower than the scores of three of the individual
streams.

The audio branch achieves a ROC-AUC of 0.5054 [0.4747, 0.5371], which is
consistent with chance-level performance. The sync branch's anomaly score
remains below 0.5 for all three seeds. Taken together, these results suggest
that the branches do not contribute equally to cross-corpus detection: the
visual branch transfers more successfully, while the audio and sync branches
do not provide a consistent improvement.

\subsection{Seed spread exceeds the measured fusion differences}

\begin{table}[htbp]
\caption{DFDC ROC-AUC by training seed}
\label{tab:seeds}
\centering
\begin{tabular}{lrrrr}
\toprule
System & 17 & 29 & 43 & Spread \\
\midrule
visual branch & 0.7611 & 0.6963 & 0.6878 & 0.0733 \\
fused system & 0.7500 & 0.6698 & 0.6422 & \textbf{0.1078} \\
\bottomrule
\end{tabular}
\end{table}

The three seeds produce very similar results within the training domain, with
scores agreeing to three or four decimal places. The cross-corpus results are
much less stable. As shown in Table~\ref{tab:seeds}, the fused system has a
DFDC ROC-AUC spread of 0.1078 across the three seeds. This is approximately
thirteen times the magnitude of the seed-17 fusion difference reported in
Table~\ref{tab:fusion}.

Rather than reducing variation, combining the branches appears to amplify
the effect of the training seed in this experiment. This makes it difficult
to interpret a result from a single run as representative of the system's
general performance.

We also evaluated an ensemble formed by averaging the six final visual
checkpoints, comprising EfficientNet-B0 and ConvNeXt-Tiny~\cite{convnext},
each trained with three seeds and assigned equal weights. The ensemble
achieves 0.7739 [0.7401, 0.8064] on DFDC. This is
$+0.0469$ [$+0.0286$, $+0.0666$] above the EfficientNet seed mean and
matches the best individual seed. On Celeb-DF, its performance is comparable
to that of every individual member.

There is an important qualification to this result. We selected this
particular ensemble grouping after comparing three groupings on DFDC.
Therefore, its improvement on DFDC still needs to be confirmed on a test
partition that has not been used during model selection.

\subsection{Choices that looked right on validation}

\textit{Selection on loss.}
One of the visual streams reached its lowest validation loss at epoch 1,
while its backbone was still frozen. However, the resulting checkpoint
achieved only 0.4668 ROC-AUC in-domain, which is below chance. This result
shows that selecting a checkpoint using validation loss does not necessarily
produce the best ranking performance.

Binary cross-entropy penalizes confident incorrect predictions heavily.
Consequently, a model with poor ranking ability can obtain a lower loss than
a model that ranks examples well but makes overconfident mistakes. Following
this observation, our trainers select checkpoints using validation ROC-AUC
rather than validation loss.

\textit{Frozen BatchNorm.}
Freezing the BatchNorm statistics during fine-tuning increased validation
performance from 0.5154 to 0.9058. Initially, we interpreted this as a clear
improvement. The cross-corpus comparison, however, told a different story.
Relative to frozen statistics, keeping the statistics live reduced in-domain
ROC-AUC by 0.5430 but improved DFDC performance by 0.1276.

The two configurations had been trained for different numbers of epochs,
which meant that the original comparison did not isolate the effect of
BatchNorm. After fixing the training budget, the DFDC difference decreased
to 0.0177, approximately one seventh of the original value. Live statistics
appear to act as an unintended domain adapter~\cite{adabn}, and the
in-domain performance cost is the part of the effect that remains evident
under the controlled comparison.

\subsection{The operating point fails before the ranking does}

\begin{table}[htbp]
\caption{Visual branch under corruption of the cached view}
\label{tab:corruption}
\centering
\begin{tabular}{lrrr}
\toprule
Condition & ROC-AUC & Genuine called fake & Fakes missed \\
\midrule
clean & 0.9988 & 0 of 150 & 3 of 1,498 \\
JPEG 75 & 0.9934 & \textbf{76 of 150} & 1 \\
JPEG 50 & 0.9187 & \textbf{150 of 150} & 0 \\
\bottomrule
\end{tabular}
\end{table}

The corruption experiment illustrates why ROC-AUC alone is not enough to
describe how a detector behaves at a fixed decision threshold. Under JPEG
quality 75, ROC-AUC decreases by only 0.0054, which might appear to indicate
that the model remains robust. Yet the number of genuine clips classified as
fake rises from zero to 76 out of 150, corresponding to a false alarm rate of
approximately 51 percent. At JPEG quality 50, every genuine clip in the
evaluation is classified as fake, despite the model missing none of the fake
clips.

Audio corruption produces a different failure pattern. Adding noise at
20~dB SNR reduces ROC-AUC by 0.0227, but the detector misses 857 of 859
fake clips. These results show that relatively small changes in a ranking
metric can accompany substantial changes in the decisions made at the
deployed operating point.

The model also becomes more confident when it is wrong on external datasets.
On DFDC, the visual branch misclassifies 1,942 of 2,410 clips at a threshold
of 0.5. Its mean confidence is 0.9244 on incorrect predictions, compared
with 0.8517 on correct predictions. No confidence threshold removes errors
faster than it removes correct predictions. On Celeb-DF, the same threshold
classifies 141 of 155 genuine videos as fake.

We also tested an approach that estimates how familiar a clip looks to the
model. A Gaussian distribution is fitted to frozen DINOv3 features extracted
from genuine training clips, and new clips are scored using Mahalanobis
distance. This model does not train on fake examples; instead, it measures
whether a clip differs from the distribution of genuine training footage.

At the deployed setting, the method abstains on 8.7 percent of held-out
genuine FakeAVCeleb clips. On genuine Celeb-DF clips, it intercepts
74.4 percent, with the result ranging from 49.4 to 83.1 percent across
twenty reference splits. This reduces false accusations on unfamiliar
footage from roughly 91 percent to somewhere between a quarter and a third.

\subsection{Open-world generated video}

Before reporting the DF26 results, we identified and removed four potential
shortcuts: scene content, real clips appearing in multiple splits, sibling
modes of the same generator, and the presence of audio. The audio difference
was particularly important because it could predict the label by itself:
every commercial clip contains an audio track, while almost none of the
other clips do. We also found that the real clips were the only ones recorded
at 25 or 29.97 frames per second. To reduce this difference, we re-encoded
all clips to $640\times360$ at 24 frames per second.

\begin{table}[htbp]
\caption{Unseen-generator ROC-AUC on DF26, frozen encoder and linear probe}
\label{tab:df26}
\centering
\begin{tabular}{lrl}
\toprule
Encoder & Macro & Paired against DINOv3 \\
\midrule
Perception Encoder~\cite{pe} & \textbf{0.7527} & $+0.0741$ [$+0.0121$, $+0.1317$] \\
CLIP ViT-B/16~\cite{clip} & 0.7454 & $+0.0668$ [$+0.0049$, $+0.1289$] \\
DINOv3 ViT-S/16 & 0.6786 & reference \\
\bottomrule
\end{tabular}
\end{table}

Perception Encoder achieves the highest macro ROC-AUC across the seven
held-out generators, as shown in Table~\ref{tab:df26}. Its scores range from
0.8496 on Veo~3.1 to 0.6542 on LTX~2.3. However, these results require an
important qualification. All three encoders received ImageNet-normalized
pixels. This preprocessing is appropriate for DINOv3 but not for the other
two encoders, meaning that the comparison was conducted under mismatched
normalization conditions. The results should therefore not be treated as
lower bounds. We consider the comparison preliminary until the features are
re-extracted using the normalization statistics specified for each encoder.

The results also show that the choice of region of interest matters. When
we evaluate a full-frame probe on FaceForensics++, holding out one
manipulation family in each arm, the macro score is close to chance at
0.5387. Using face crops with the fine-tuned visual stream, retrained with
one family held out each time, reduces the macro score by 0.1231 compared
with a model trained on all six families. The graphics-based FaceSwap method
achieves 0.5195.

These findings suggest that full frames may be better suited to detecting
fully generated video, whereas face crops are more appropriate when the
task involves manipulation of a real face. The distinction is relevant
because the visual evidence available to a detector differs between these
two types of forgery.

\subsection{Encoders the cross-modal design named}

The original stream design specified AV-HuBERT~\cite{avhubert} for lip-sync
detection and HSEmotion~\cite{hsemotion} for the face component of the
emotion stream. Neither could be used directly in its released form in our
environment. AV-HuBERT depends on fairseq, which requires an older PyTorch
version, while the HSEmotion model is distributed as a pickled model that
the current version of timm cannot load.

To work around these compatibility issues, we reimplemented AV-HuBERT Base
in plain PyTorch and loaded the HSEmotion weights into the same
EfficientNet-B0 architecture. Both implementations were checked against
the official code using identical inputs, and their features and logits
agree to within $2\times10^{-6}$. AV-HuBERT is run separately on the lips
and the audio because its pretraining maps both modalities into a shared
representation space.

On 98 percent of 128 genuine training clips, correctly aligned audio matches
the lips better than audio shifted by half a second. A lip-sync stream using
these features achieves a validation ROC-AUC of 0.9416, compared with
0.7593 for the Wav2Vec2-based stream. The selected checkpoint comes from
epoch 2 of a five-epoch run, before the encoder is unfrozen at epoch 3.

Despite the improvement in validation performance, the diagonal attention
mass remains between 0.05875 and 0.05919, close to the uniform value of
0.0592. This means that the validation result does not establish that the
trained attention head has learned temporal alignment. The shifted-audio
experiment provides evidence about the pretrained features, but it does not
confirm that the downstream detector uses those features through the intended
attention mechanism.

These results are preliminary and come from seed 17 only. Held-out
FakeAVCeleb, DFDC, and fusion scores for the replacement encoders have not
yet been computed. The HSEmotion training run has also not produced a
completed result for this comparison. Validation performance alone is
therefore insufficient to conclude that either replacement improves
cross-corpus deepfake detection.

\section{Limitations}

Several limitations affect how the results should be interpreted. First,
all Design B streams were trained with only one seed, although the protocol
specifies three. As Table~\ref{tab:seeds} shows, the choice of seed can
substantially change cross-corpus performance. Both Design A and Design B
were trained on FakeAVCeleb, so their results are comparable to each other
but should not be directly compared with published results obtained under
different protocols. Only the FaceForensics++ experiment follows the
published cross-dataset protocol, and that experiment also uses a single
seed.

None of the training data was collected in the wild. More importantly, the
cross-modal streams do not provide evidence that they have learned the
correspondence their design is intended to capture. One of the original
research exit criteria required the confidence interval of the synchronization
offset to exclude chance. This was never measured because the branch's
training history recorded only a combined loss.

The first track also has limited coverage of fully generated video. It
detects none of the 10 MNW clips from one talking-head generator. The second
track relies on a single corpus and a single split, and DF26 is not suitable
for audio experiments. These limitations mean that the findings should be
interpreted within the evaluation settings described here. Additional seeds,
held-out tests, and direct measurements of temporal correspondence are needed
before drawing broader conclusions about the generalization of the proposed
systems.

\section{Conclusion}

Our experiments show that cue-specific fusion did not outperform the visual
baseline across corpora. Combining every stream also performed worse than
three of the individual streams. Under the evaluation conditions used in this
study, the original hypothesis was therefore not supported.

The experiments also revealed why performance scores need to be examined
alongside checks of model behaviour. One stream achieved a ROC-AUC of 0.9991
without showing meaningful changes in its attention distribution. The visual
model was sensitive to codec differences, freezing BatchNorm statistics
changed the apparent balance between in-domain and cross-corpus performance,
and selecting a checkpoint by validation loss produced below-chance
performance. We also found that the training seed alone could shift the
cross-corpus ROC-AUC by 0.1078.

These findings were not apparent from the training objective or the headline
performance metrics alone. They emerged from measurements that we recorded
without directly optimizing them. The experiments also helped us correct our
own interpretation of an earlier result: an apparent effect became roughly
seven times smaller after the training budget was controlled.

Taken together, the results suggest that reporting a single cross-corpus
score to four decimal places can give a misleading impression of certainty
when the choice of seed can produce much larger differences. A reliable
evaluation should examine not only whether a detector performs well, but
also whether its internal behaviour supports the explanation for its
predictions. Reporting seed variation, checking the mechanisms used by
individual streams, and testing for shortcuts are therefore important parts
of evaluating audio-visual deepfake detectors.

\begin{thebibliography}{00}
\bibitem{fakeavceleb} H. Khalid, S. Tariq, M. Kim, and S. S. Woo,
``FakeAVCeleb: A novel audio-video multimodal deepfake dataset,'' in
\textit{Proc. NeurIPS Datasets and Benchmarks Track}, 2021.
\bibitem{lavdf} Z. Cai, K. Stefanov, A. Dhall, and M. Hayat, ``Do you really
mean that? Content driven audio-visual deepfake dataset and multimodal method
for temporal forgery localization,'' in \textit{Proc. DICTA}, 2022.
\bibitem{celebdf} Y. Li, X. Yang, P. Sun, H. Qi, and S. Lyu, ``Celeb-DF: A
large-scale challenging dataset for DeepFake forensics,'' in \textit{Proc.
CVPR}, 2020.
\bibitem{dfdc} B. Dolhansky, J. Bitton, B. Pflaum, J. Lu, R. Howes, M. Wang,
and C. Canton Ferrer, ``The DeepFake Detection Challenge (DFDC) dataset,''
arXiv:2006.07397, 2020.
\bibitem{ffpp} A. R\"ossler, D. Cozzolino, L. Verdoliva, C. Riess, J. Thies,
and M. Nie{\ss}ner, ``FaceForensics++: Learning to detect manipulated facial
images,'' in \textit{Proc. ICCV}, 2019.
\bibitem{syncnet} J. S. Chung and A. Zisserman, ``Out of time: Automated lip
sync in the wild,'' in \textit{Proc. ACCV Workshops}, 2016.
\bibitem{avhubert} B. Shi, W.-N. Hsu, K. Lakhotia, and A. Mohamed, ``Learning
audio-visual speech representation by masked multimodal cluster prediction,''
in \textit{Proc. ICLR}, 2022.
\bibitem{sbi} K. Shiohara and T. Yamasaki, ``Detecting deepfakes with
self-blended images,'' in \textit{Proc. CVPR}, 2022.
\bibitem{univfd} U. Ojha, Y. Li, and Y. J. Lee, ``Towards universal fake image
detectors that generalize across generative models,'' in \textit{Proc. CVPR},
2023.
\bibitem{shortcut} R. Geirhos, J.-H. Jacobsen, C. Michaelis, R. Zemel,
W. Brendel, M. Bethge, and F. A. Wichmann, ``Shortcut learning in deep neural
networks,'' \textit{Nature Machine Intelligence}, vol. 2, pp. 665-673, 2020.
\bibitem{audit} ``What do deepfake benchmarks measure? An audit using frozen
self-supervised representations,'' arXiv:2606.26384, 2026.
\bibitem{efficientnet} M. Tan and Q. V. Le, ``EfficientNet: Rethinking model
scaling for convolutional neural networks,'' in \textit{Proc. ICML}, 2019.
\bibitem{mtcnn} K. Zhang, Z. Zhang, Z. Li, and Y. Qiao, ``Joint face detection
and alignment using multitask cascaded convolutional networks,'' \textit{IEEE
Signal Process. Lett.}, vol. 23, no. 10, pp. 1499-1503, 2016.
\bibitem{wav2vec2} A. Baevski, Y. Zhou, A. Mohamed, and M. Auli, ``wav2vec 2.0:
A framework for self-supervised learning of speech representations,'' in
\textit{Proc. NeurIPS}, 2020.
\bibitem{dinov3} O. Sim\'eoni \textit{et al.}, ``DINOv3,'' arXiv:2508.10104,
2025.
\bibitem{mnw} Microsoft, Northwestern University, and WITNESS, ``MNW deepfake
detection benchmark,'' \url{https://github.com/microsoft/MNW}.
\bibitem{df26} ``DF26,'' dataset,
\url{https://huggingface.co/datasets/DF26/DF26}, 2026.
\bibitem{bootstrap} B. Efron and R. J. Tibshirani, \textit{An Introduction to
the Bootstrap}. New York: Chapman \& Hall, 1993.
\bibitem{convnext} Z. Liu, H. Mao, C.-Y. Wu, C. Feichtenhofer, T. Darrell, and
S. Xie, ``A ConvNet for the 2020s,'' in \textit{Proc. CVPR}, 2022.
\bibitem{adabn} Y. Li, N. Wang, J. Shi, X. Hou, and J. Liu, ``Adaptive batch
normalization for practical domain adaptation,'' \textit{Pattern Recognition},
vol. 80, pp. 109-117, 2018.
\bibitem{pe} D. Bolya \textit{et al.}, ``Perception Encoder: The best visual
embeddings are not at the output of the network,'' arXiv:2504.13181, 2025.
\bibitem{clip} A. Radford \textit{et al.}, ``Learning transferable visual
models from natural language supervision,'' in \textit{Proc. ICML}, 2021.
\bibitem{hsemotion} A. V. Savchenko, ``Facial expression and attributes
recognition based on multi-task learning of lightweight neural networks,'' in
\textit{Proc. IEEE SISY}, 2021.
\end{thebibliography}

\end{document}